\section{Алгоритмы распознавания языка}
\label{sec:algorithms}

По варианту~13 нужно понять, на каких языках написан PDF: на русском, на немецком или на обоих сразу. Для этого сравниваются три способа \mdash короткие слова, частотные слова и нейросеть.

Сначала по учебным текстам для каждого языка собирается поисковый образ языка (ПОЯ): какие слова или куски текста для него обычны. У нового документа по тем же правилам строится поисковый образ документа (ПОД). Дальше для каждого языка считается оценка принадлежности. Для слов это произведение вероятностей, для нейросети \mdash доля кириллических и латинских N-грамм. Язык считается найденным, если его доля не меньше~$0{,}15$. Поэтому смешанный документ показывает оба языка, а не только тот, которого чуть больше. Дополнительно считается расстояние OoP (Out-of-Place): насколько по-разному выстроены списки частых элементов.

Учебные тексты лежат в каталоге \texttt{corpus/train}: по семь медицинских заметок на каждый язык, примерно 50~Кбайт русского текста и 85~Кбайт немецкого. Для проверки отдельно приготовлены шестнадцать текстов: обычные фразы, длинные слова из частотного списка, редкие длинные слова и смешанные русско-немецкие фрагменты. В обучение они не входят.

Перед сравнением текст приводится к нижнему регистру, буква <<ё>> заменяется на <<е>>. Немецкие умлауты сохраняются. Русские слова приводятся к начальной форме, немецкие только к нижнему регистру. Короткие служебные слова вроде <<и>> или <<der>> не выбрасываются: по ним языки хорошо отличаются.

В методе коротких слов в образ языка попадают слова не длиннее пяти букв, если в учебных текстах они встретились больше трёх раз. Чем слово чаще, тем выше его вероятность. Если слова в образе нет, ему даётся очень маленькая вероятность, чтобы расчёт не обрывался. Документ оценивается как произведение вероятностей его коротких слов:

\begin{equation}
	P(\mathrm{lang}\mid\mathrm{doc}) = \prod_{l \in L_{\mathrm{short}}} P(l\mid\mathrm{lang}).
	\label{eq:short}
\end{equation}

Если коротких слов в тексте нет, метод язык не назначает. Иначе берутся все языки, чья доля после softmax не ниже порога. Считают сумму логарифмов: само произведение на длинном тексте слишком быстро становится нулём.

В методе частотных слов в образ входят пятьдесят самых обычных слов, длина здесь не важна. С документом сравниваются только слова из этого короткого списка. Формула та же, что и~\eqref{eq:short}. Если ни одно слово документа не входит в список, метод тоже ничего не обнаруживает.

Нейросеть смотрит не на целые слова, а на короткие куски текста длиной от одного до пяти символов: такой кусок в тексте есть или нет. По доли кириллицы и латиницы она показывает оба языка, если оба алфавита встречаются в документе. Расстояние OoP для неё не считается.

OoP сравнивает два списка, упорядоченных от частых элементов к редким. Для каждого элемента смотрят, насколько разъехались его места в списках. Если элемента во втором списке нет, разница берётся большой:

\begin{equation}
	d(A,B) = \sum_{g \in A \cup B} \lvert \mathrm{pos}_A(g) - \mathrm{pos}_B(g) \rvert.
	\label{eq:oop}
\end{equation}

Число растёт вместе с длиной списков, поэтому у коротких слов и у пятидесяти частых слов оно разного масштаба. Решение словесные методы принимают по вероятности~\eqref{eq:short}, а не по минимуму OoP.
